\begin{tikzpicture}[font=\figurefont\small,
  state/.style={draw=BookRule,fill=FigureInput!12,align=center,
    minimum width=18mm,minimum height=8mm},
  arr/.style={-{Stealth[length=2mm]},draw=BookInk,line width=.8pt}]
  \node[anchor=west] at (-1,3.6) {(a) 进展状态可观察，回合起点可重置};
  \foreach \i/\x/\txt in {0/0/$s_0$,1/3/$s_1$,2/6/$s_2$,3/9/$\cdots$,4/12/$s_H$}
    \node[state] (v\i) at (\x,2.5) {\txt};
  \foreach \i/\j/\a in {0/1/0,1/2/1,2/3/2,3/4/{H-1}}
    \draw[arr] (v\i) -- node[above] {$a^*_{\a}$} (v\j);
  \node[state,fill=FigureLoss!12] (fail) at (6,.3) {可观察失败态};
  \draw[arr,BookAmber] (v0.south) -- (0,.3) -- (fail.west);
  \draw[arr,BookAmber] (v1.south) -- (3,1.3) -- (6,1.3) -- (fail.north);
  \node[anchor=west] at (7.7,.3) {其余层错误动作同样失败};
  \node[anchor=west] at (0,1) {错误动作};
  \node[anchor=west] at (-1,-1.3) {(b) 隐藏进展，只在末端观察整串是否正确};
  \foreach \i/\x/\txt in {0/0/时刻0,1/3/时刻1,2/6/时刻2,3/9/$\cdots$,4/12/末端反馈}
    \node[state,fill=BookPaper] (h\i) at (\x,-2.4) {\txt};
  \foreach \i/\j in {0/1,1/2,2/3,3/4}
    \draw[arr] (h\i) -- (h\j);
  \node[align=center,text width=138mm] at (6,-4)
    {上：可逐层辨识并复用正确前缀。下：一次失败只排除一个完整动作串。\\若改变可观察信息或加入任意状态查询，困难程度也随之改变。};
\end{tikzpicture}
